% GPT-6 Astra Pro; effort not exposed; /root/replacement; 22 November 2026.
\section[HD-F1-R12: the two-slope whole-character obstruction]{HD-F1-R12: two exponential slopes cannot give the balanced whole character}
\label{sec:followup-rm-two-slope}

Retain the fourfold and balanced character of R07. There is no
actual perfect object $E$ with both
\[
 \operatorname{ch}(E)=N e^{-B}(1-qA_1)(1-qA_2),\qquad
 N\ne0,\quad q\ne0,\quad A_i=\theta_i^2/2,
 \tag{F1R12.1}
\]
and a character expressible using at most two exponential slopes:
\[
             \operatorname{ch}(E)=a e^{\ell_1}+b e^{\ell_2},
 \qquad a,b\in\mathbf Q,\quad
             \ell_1,\ell_2\in\operatorname{NS}(X)_{\mathbf Q}.
 \tag{F1R12.2}
\]
The slopes need not lie in $F\theta$. This obstruction holds
before imposing any finite action or semiregularity condition.
In particular it excludes every actual adjacent two-cohomology
object whose two coherent cohomology sheaves have nonzero-rank
exponential characters, including every actual coupling class.
It also excludes all actual complexes whose coherent cohomology
sheaves have nonzero-rank exponential characters at at most two
slopes, however many cohomology degrees occur. The same statement
holds for actual finite filtrations by distinguished triangles
with such factors. It is the character relation of those actual
triangles, rather than a splitting assertion, that is used below.

\paragraph{The exact exterior-square map.}
Write $V=H^0(T_X)$ and $W=H^1(\mathcal O_X)$, both of
dimension four. For a divisor class $d\in H^1(\Omega_X^1)$,
let
\[
 L_d:V\longrightarrow W,\qquad v\longmapsto\iota_vd.
 \tag{F1R12.3}
\]
For $\beta\in H^2(\Omega_X^2)$ define the linear map
\[
 R_\beta:\bigwedge^2V\longrightarrow\bigwedge^2W,
 \qquad R_\beta(v\wedge w)=-\iota_w\iota_v\beta.
 \tag{F1R12.4}
\]
This sign fixes the exact identification: contraction is a
graded derivation, $\iota_w(\iota_vd)=0$, and therefore
\[
 -\iota_w\iota_v(d^2/2)
       =(\iota_vd)\wedge(\iota_wd),\qquad
                  R_{d^2/2}=\bigwedge^2L_d.
 \tag{F1R12.5}
\]
For any linear map $L$ of rank $s$, its image on the exterior
square is exactly $\bigwedge^2\operatorname{im}L$. Thus
\[
 \operatorname{rank}R_{d^2/2}=\binom{s}{2}
                    \in\{0,1,3,6\},\qquad 0\le s\le4.
 \tag{F1R12.6}
\]
There is no diagonalizability or positivity assumption in this
rank identity.

For $C=A_1+A_2$, the actual eigenspace decompositions are
$V=V_1\oplus V_2$ and $W=W_1\oplus W_2$, each block
of dimension two. Compatibility and nondegeneracy of the
polarization make $L_{\theta_i}:V_i\to W_i$ an isomorphism,
with zero restriction to the other block. Formula (F1R12.5)
therefore gives
\[
 R_C=\bigwedge^2L_{\theta_1}
                   +\bigwedge^2L_{\theta_2},\qquad
                         \operatorname{rank}R_C=2.
 \tag{F1R12.7}
\]
The two images are the independent lines
$\bigwedge^2W_1$ and $\bigwedge^2W_2$; the mixed
domain $V_1\wedge V_2$ is killed. Consequently no nonzero
scalar multiple of $C$ is the square of a divisor class.
This is an explicit map obstruction in the whole $H^{2,2}$
space, without restricting the divisor to the RM span.

\paragraph{Centering the two actual character terms.}
The degree-zero and degree-one parts of (F1R12.1)--(F1R12.2)
give
\[
       a+b=N,\qquad a\ell_1+b\ell_2=-NB.
 \tag{F1R12.8}
\]
If one coefficient vanishes, multiplying by $e^B$ makes the
single remaining exponential equal to $N$ in degree zero
and have zero slope. Its degree-two component then vanishes,
contrary to $-NqC\ne0$. Thus both coefficients are nonzero.
Put $d=\ell_1-\ell_2$. The precise centered slopes are
\[
                B+\ell_1=\frac bN d,\qquad
                B+\ell_2=-\frac aN d.
 \tag{F1R12.9}
\]
The degree-two identity is now
\[
 \frac12\left(a(B+\ell_1)^2+b(B+\ell_2)^2\right)
       =\frac{ab}{N}\frac{d^2}{2}=-NqC.
 \tag{F1R12.10}
\]
Applying the actual map (F1R12.4) to this equality forces
$\bigwedge^2L_d$ to have rank two, contradicting
(F1R12.6). This proves the exclusion.

\paragraph{Actual adjacent objects and framework scope.}
For an adjacent object with $H^0(E)=A$, $H^1(E)=T$, the
canonical triangle is $A\to E\to T[-1]\to A[1]$ with
connecting morphism $\delta\in\operatorname{Ext}^2(T,A)$.
For exponential endpoint characters
$\operatorname{ch}(A)=r_Ae^{\ell_A}$ and
$\operatorname{ch}(T)=r_Te^{\ell_T}$, its actual
Grothendieck relation gives (F1R12.2) with
$a=r_A$, $b=-r_T$. The proof covers every $\delta$ and
does not use the dimension of $\operatorname{Hom}(A,T)$.
For several cohomology sheaves, the canonical truncation
triangles give $[E]=\sum_j(-1)^j[H^j(E)]$; aggregating
their actual ranks at either slope gives the same $a,b$.
For a finite triangle filtration, repeated additivity gives
the identical character equation.

The character is the GRR realization of these exact $K_0$
relations. Adams operations multiply its codimension-two
part by the square of the operation index, preserving the
rank obstruction in (F1R12.10); the scalar is explicit and
nonzero. Simple semihomogeneous bundles supply examples of
the exponential endpoint condition by Mukai's actual
isogeny-line construction, with the full character derived
in R09 from its pullback decomposition. No classification of
all coherent sheaves is being asserted. The proof supplies
no conclusion for two arbitrary full-rank characters, for
three or more distinct slopes, or for a transformed
rank-zero character not satisfying (F1R12.2). The source
comparison uses Mukai's author text, equation (4.8), and
the published Orlov Proposition 4.5 already checked in R09;
the exterior-square calculation above is proved in full.
